% Generated draft: requires mathematical and source review.
\documentclass[11pt]{article}
\usepackage[margin=25mm]{geometry}
\usepackage{fontspec}
\setmainfont{Latin Modern Roman}
\usepackage{amsmath,amssymb,mathtools,unicode-math}
\setmathfont{Latin Modern Math}
\usepackage{enumitem,needspace,xurl,hyperref}
\setlength{\parindent}{0pt}
\setlength{\parskip}{5pt}
\setlength{\emergencystretch}{3em}
\newcommand{\N}{\mathbb{N}}
\newcommand{\Z}{\mathbb{Z}}
\newcommand{\Q}{\mathbb{Q}}
\newcommand{\R}{\mathbb{R}}
\newcommand{\C}{\mathbb{C}}
\newcommand{\F}{\mathbb{F}}
\newcommand{\A}{\mathbb{A}}
\newcommand{\PP}{\mathbb{P}}
\newcommand{\E}{\mathbb{E}}
\newcommand{\eps}{\varepsilon}
\newcommand{\dd}{\,\mathrm{d}}
\begin{document}
\section*{UnsolvedMath ID: 30004829}
\textbf{Source problem code:} catalogue entry
\subsection{Definitions and mathematical statement}
Let \(G\) be a finitely generated group and let \(\C G\) be its group algebra. A left-left Yetter--Drinfeld module over \(\C G\) is a complex vector space \(V\) equipped with a left \(G\)-action and a grading
\[
V=\bigoplus_{g\in G}V_g,\qquad h\cdot V_g=V_{hgh^{-1}}\quad(g,h\in G).
\]
Its support \(\{g\in G:V_g\ne0\}\) is therefore a union of conjugacy classes. Write \({}^{\C G}_{\C G}\mathcal{YD}\) for this category. Its braiding is
\[
c(v\otimes w)=g\cdot w\otimes v\qquad(v\in V_g,\ w\in V).
\]
No finite-dimensionality assumption on \(V\) is stated in the supplied problem.

Give the tensor algebra \(T(V)=\bigoplus_{n\ge0}V^{\otimes n}\) its braided graded Hopf algebra structure with \(V\) primitive. The Nichols algebra is
\[
\mathcal B(V)=T(V)/I(V),
\]
where \(I(V)\) is the largest homogeneous braided Hopf ideal contained in \(\bigoplus_{n\ge2}V^{\otimes n}\). The bosonization \(\mathcal B(V)\#\C G\) has underlying vector space \(\mathcal B(V)\otimes\C G\) and multiplication
\[
(b\#g)(b'\#h)=b(g\cdot b')\#gh.
\]
For a unital complex algebra \(A\), use the definition
\[
\operatorname{GKdim}A=
\sup_{\substack{W\subseteq A\text{ finite-dimensional}\\1\in W}}
\limsup_{n\to\infty}\frac{\log\dim_{\C}(W^n)}{\log n},
\]
where \(W^n\) is the linear span of products of \(n\) elements of \(W\). This definition also applies when finite generation of \(A\) has not been assumed.

\textbf{Classification problem (source Question 2).} For each finitely generated group \(G\), classify the objects \(V\in{}^{\C G}_{\C G}\mathcal{YD}\) satisfying
\[
\operatorname{GKdim}\bigl(\mathcal B(V)\#\C G\bigr)<\infty.
\]
The usual interpretation is classification up to linear isomorphisms preserving both the \(G\)-action and the \(G\)-grading, with \(G\) fixed; the supplied source does not expressly specify the equivalence relation.

\textbf{Scope and source qualifications.} The surrounding report concerns pointed Hopf algebras over \(\C\) of finite Gelfand--Kirillov dimension. It states, invoking Gromov, that in its finite-growth setting one may assume \(G\) is nilpotent-by-finite, meaning that \(G\) has a nilpotent subgroup of finite index. The supplied statement itself quantifies over all finitely generated groups, whereas the supplied literature assessment refers to a ``nilpotent-group setting.'' These descriptions are retained as distinct source qualifications; no additional nilpotence hypothesis is imposed here.

The original excerpt presents Question 2 as the core of Question 1, which asks to classify, or at least characterize, coradically graded connected Hopf algebras \(E=\bigoplus_{n\ge0}E_n\) in this category with \(\operatorname{GKdim}(E\#\C G)<\infty\). It describes \(E\), with \(V=E_1\), as a post-Nichols algebra and identifies the subalgebra generated by \(V\) with \(\mathcal B(V)\). Dually, it calls a graded connected Hopf algebra \(R\) generated by \(R_1\simeq V\) a pre-Nichols algebra. The excerpt's parenthetical assertion about the ``only pre- and post-Nichols algebra'' is ambiguous and is not used as a hypothesis or theorem in this draft.
\subsection{Short English statement}
For each finitely generated group \(G\), classify the complex Yetter--Drinfeld modules whose Nichols algebra, after bosonization with \(\C G\), has finite Gelfand--Kirillov dimension. The source's contextual restriction to nilpotent-by-finite groups and its literature assessment's reference to nilpotent groups require reconciliation. The supplied partially solved label is dated curation, not an independently verified mathematical conclusion.
\subsection{Sources}
\begin{itemize}
\item Supplied local UnsolvedMath registry record: permanent ID \(30004829\), title ``Finite-Growth Nichols Algebras over Finitely Generated Groups,'' exact source problem code \texttt{catalogue entry}, Oberwolfach Reports Open Problems collection, source Question 2. The record gives proposed year \(2021\) and supplies no named proposer.
\item Supplied report citation: \emph{Quantum Groups -- Algebra, Analysis and Category Theory} (2022), Oberwolfach Reports, pp.~2397--2458, DOI \texttt{10.4171/owr/2021/44}. \url{https://doi.org/10.4171/owr/2021/44}. This is also the literature source labelled ``Oberwolfach Report 44.'' The citation's publication year \(2022\) and the record's proposed year \(2021\) are preserved separately.
\item Supplied literature source labelled ``Angiono--Garcia Iglesias'': \url{https://arxiv.org/abs/2212.08169}.
\item Supplied literature source labelled ``Nichols algebras over Q8'': \url{https://arxiv.org/abs/2312.10505}.
\item Literature assessment dated exactly \texttt{2026-08-22}: the supplied curation labels the problem partially solved and states that finite-GK Nichols algebras of diagonal type are characterized by finite root systems, with further results for selected groups, while the full classification in its stated nilpotent-group setting remains incomplete. This is supplied dated curation, not independently verified here; its time-sensitive status and cited claims require independent checking.
\item The registry labels the statement \texttt{corrected\_verified} and reports that the linked report was inspected and its numbered Question 2 transcribed as a self-contained statement. These are supplied source-verification claims, not verification performed for this draft. The registry's supplied update timestamp is \texttt{2026-09-16T12:06:53Z}.
\item The original excerpt attributes its classification method to H.-J. Schneider and an unnamed speaker, invokes a result of Gromov for the group-growth reduction, and refers to a survey as reference [1]. The supplied record does not identify the speaker or provide the bibliographic entry for [1].
\end{itemize}
\end{document}
